\section{Harness Engineering: cómo entra el modelo en un flujo de trabajo real}

En la entrevista se menciona Harness Engineering: además del modelo, hace falta una capa de ingeniería que conecte el modelo con el flujo de trabajo real. Sin harness, el modelo es solo una interfaz potente; con harness, el modelo puede conectarse con el repositorio de código, las pruebas, los permisos, los registros, los reintentos, el despliegue y la retroalimentación.

\begin{figure}[H]
\centering
\includegraphics[width=0.92\textwidth]{harness-engineering.png}
\caption{Harness Engineering: la capa intermedia entre el modelo y el flujo de trabajo real.}
\end{figure}

\begin{knowledgebox}{Lectura de la figura: por qué no basta con que el modelo sea potente}
En la figura, el modelo es solo la primera capa. El harness se encarga de los prompts, las herramientas, los permisos, el estado y los reintentos; Workflow es el repositorio de código real y los procesos de negocio; Feedback devuelve al sistema las pruebas, la retroalimentación de los usuarios, los costos y los casos de fallo; solo entonces Improvement puede pasar al posentrenamiento y a la iteración del producto.
\end{knowledgebox}

\subsection{Términos clave: vocabulario relacionado con harness}

\noindent\begin{tabularx}{\textwidth}{p{0.22\textwidth}p{0.32\textwidth}X}
\toprule
Término & Problema que resuelve & Significado en este episodio \\
\midrule
Harness & Sistema externo que conecta el modelo con herramientas y flujos de trabajo & Determina si el modelo puede trabajar de forma estable. \\
Workflow & Flujo de tareas real, como issues, pruebas y despliegue & El valor del modelo se materializa, en última instancia, dentro del workflow. \\
Feedback & Pruebas, usuarios, costos, casos de fallo & Permite que el sistema aprenda y se mejore. \\
Permission & Permisos y límites de seguridad & Evita que el modelo realice operaciones sin autorización. \\
Retry / Recovery & Reintento y recuperación tras un error & Determina si las tareas de largo alcance pueden completarse de forma estable. \\
\bottomrule
\end{tabularx}

\subsection{Por qué se relaciona con el posentrenamiento}

El harness no es solo ingeniería de producto; también influye de vuelta en el entrenamiento. Un buen harness produce casos de fallo más claros, trayectorias de tareas más realistas y retroalimentación más verificable. Todo esto puede convertirse en datos de posentrenamiento. En el EP138, cuando Luo Fuli (罗福莉) habla del Agent post-train, enfatiza el entorno de tareas y el rollout; es la misma lógica que aquí.

\begin{warningbox}{No tratar el harness como un envoltorio de UI}
La UI es solo una pequeña parte del harness. Un harness real incluye gestión de estado, ejecución de herramientas, seguridad, registros, evaluación y retorno de la retroalimentación. Hacer solo una interfaz atractiva no permite que el modelo entre de forma confiable en producción.
\end{warningbox}

\subsection{Resumen de la sección}

Harness Engineering es el puente por el que el modelo se convierte en un sistema productivo. Conecta las capacidades del modelo con tareas reales y devuelve la retroalimentación de esas tareas al entrenamiento y a la iteración del producto.
